
\documentclass{article}
\usepackage{colortbl}
\usepackage{makecell}
\usepackage{multirow}
\usepackage{supertabular}

\begin{document}

\newcounter{utterance}

\centering \large Interaction Transcript for game `cladder', experiment `full\_v1.5\_default', episode 3456 with qwen3.5{-}4b.
\vspace{24pt}

{ \footnotesize  \setcounter{utterance}{1}
\setlength{\tabcolsep}{0pt}
\begin{supertabular}{c@{$\;$}|p{.15\linewidth}@{}p{.15\linewidth}p{.15\linewidth}p{.15\linewidth}p{.15\linewidth}p{.15\linewidth}}
    \# & \multicolumn{2}{c}{Player} && \multicolumn{2}{c}{Game Master} \\
    \hline

    \theutterance \stepcounter{utterance}  
    & & & \multicolumn{4}{p{0.6\linewidth}}{
        \cellcolor[rgb]{0.9,0.9,0.9}{
            \makecell[{{p{\linewidth}}}]{
                \texttt{\tiny{[P1$\langle$GM]}}
                \texttt{You are an expert in causal inference. The following question is not a typical commonsense query, but rather a meticulously designed question created by a professor specializing in causal inference, intended to assess the students' mastery of the course content.} \\
\texttt{Imagine a self{-}contained, hypothetical world with only the following conditions, and without any unmentioned factors or causal relationships: Solar eclipse has a direct effect on Alice arriving to school. Traffic has a direct effect on Alice arriving to school. Method 1: We look at how solar eclipse correlates with Alice arriving to school case by case according to traffic. Method 2: We look directly at how solar eclipse correlates with Alice arriving to school in general. To understand how solar eclipse affects Alice arriving to school, is it more correct to use the Method 1 than Method 2?} \\
\texttt{Start your answer with 'Yes' or 'No', followed by additional reasoning or evidence to support your explanation.} \\
            }
        }
    }
     \\ \\

    \theutterance \stepcounter{utterance}  
    & \multicolumn{4}{p{0.6\linewidth}}{
        \cellcolor[rgb]{0.561,0.737,0.561}{
            \makecell[{{p{\linewidth}}}]{
                \texttt{\tiny{[P1$\rangle$GM]}}
                \texttt{Yes} \\
\\ 
\texttt{**Reasoning:**} \\
\\ 
\texttt{In causal inference, the goal is to estimate the **causal effect** of a treatment (Solar Eclipse) on an outcome (Alice arriving to school). To do this correctly, we must account for **confounding variables**—factors that influence both the treatment and the outcome, creating a spurious correlation if ignored.} \\
\\ 
\texttt{In this hypothetical world:} \\
\texttt{1.  **Solar Eclipse** directly affects Alice's arrival.} \\
\texttt{2.  **Traffic** directly affects Alice's arrival.} \\
\texttt{3.  Crucially, the prompt implies a scenario where these factors interact or coexist. In real{-}world causal structures (and implied by the distinction between the two methods), "Traffic" acts as a **confounder** or a **collider** depending on the specific temporal or structural relationship, but most likely here it serves as a confounder if Solar Eclipses and Traffic conditions are correlated (e.g., people drive more during eclipses, or traffic patterns change due to the event).} \\
\\ 
\texttt{*   **Method 2** calculates the crude correlation between Solar Eclipse and Alice's arrival. If Traffic is a confounder (i.e., it causes both the likelihood of an eclipse{-}related behavior and the traffic conditions, or if the eclipse and traffic are correlated), Method 2 will capture the effect of the eclipse **plus** the effect of traffic. This leads to **confounding bias**, resulting in an incorrect estimate of the causal effect.} \\
\\ 
\texttt{*   **Method 1** involves looking at the correlation "case by case according to traffic." This describes **stratification** (or conditioning on the confounder). By holding Traffic constant within each stratum (e.g., comparing eclipse effects only when traffic is heavy, and only when traffic is light), we block the backdoor path through Traffic. This isolates the direct effect of the Solar Eclipse on Alice's arrival, removing the bias introduced by Traffic.} \\
\\ 
\texttt{Therefore, to understand the *true* causal effect of the solar eclipse, one must control for the confounding variable (Traffic). Method 1 achieves this control, whereas Method 2 does not. Thus, Method 1 is the correct approach.} \\
            }
        }
    }
    & & \\ \\

    \theutterance \stepcounter{utterance}  
    & & & \multicolumn{2}{p{0.3\linewidth}}{
        \cellcolor[rgb]{0.9,0.9,0.9}{
            \makecell[{{p{\linewidth}}}]{
                \texttt{\tiny{[GM$|$GM]}}
                \texttt{Yes} \\
            }
        }
    }
    & & \\ \\

    \theutterance \stepcounter{utterance}  
    & & & \multicolumn{2}{p{0.3\linewidth}}{
        \cellcolor[rgb]{0.9,0.9,0.9}{
            \makecell[{{p{\linewidth}}}]{
                \texttt{\tiny{[GM$|$GM]}}
                \texttt{no} \\
            }
        }
    }
    & & \\ \\

    \theutterance \stepcounter{utterance}  
    & & & \multicolumn{2}{p{0.3\linewidth}}{
        \cellcolor[rgb]{0.9,0.9,0.9}{
            \makecell[{{p{\linewidth}}}]{
                \texttt{\tiny{[GM$|$GM]}}
                \texttt{game\_result = LOSE} \\
            }
        }
    }
    & & \\ \\

\end{supertabular}
}

\end{document}
